% =====================================================================
%  26.2  DP on the compressed tree -- edge weights, skipped counts, the
%  fold that equals the honest O(n) DP.
%  Exports: sec:c30-dp  def:c30-edgeweight  thm:c30-fold  lst:c30-fold
%           prob:c30-spectral  lst:c30-spectral
%  Solved CF 2252F "Spectral Components" (2600, L6), live-verified against the
%  CF problem page + the Round 1115 editorial (blog/entry/155741), whose official
%  solution is "for each colour build a Virtual Tree of its M vertices in
%  O(M log M) using LCA; run a bottom-up DP on the Virtual Tree".  Official
%  3-test sample read live; test 1 replayed to "1 3 -1 -1 -1" here.
% =====================================================================
\section{Running a DP on the Compressed Tree}\label{sec:c30-dp}

\begin{previouslybox}[Ch.~13 taught the fold; here it runs on a tree that shrank]
Chapter~13 (\(\thmref{c10-subtree-template}\)) established the bottom-up subtree
fold: process a node's children first, combine their answers into the node's.
Chapter~16 proved \(\dist(u,v)\) is pure \(\depth\) arithmetic
(\(\thmref{c11-distf}\)).  This section marries the two: run Ch.~13's fold on the
virtual tree of \(\secref{c30-construction}\) instead of the real tree, and charge
each \term{compressed edge} the exact quantity the chain it swallowed would have
contributed.  The payoff: a computation that would be \(O(n)\) per marked set
becomes \(O(k)\) per marked set -- the reason a per-colour or per-query problem that
is \(O(n\cdot\#\text{sets})\) naively drops to \(O(\sum k)\).
\end{previouslybox}

\begin{topicintro}
A virtual tree is not just fewer nodes -- it is a \emph{faithful} tree, in that any
distance, ancestor relation or connected-subgraph among the kept nodes has the same
value as in the real tree, provided every compressed edge carries its real length.
Once you accept that, the whole Chapter~13 DP toolkit transfers verbatim: you write
\(\mathit{dp}[v]=f(\mathit{dp}[\text{virtual children}])\), only each child edge also
multiplies in its weight and in the number of \emph{skipped} real nodes hanging below
it.  This section states the two carrying rules precisely and spends them on CF 2252F,
where each colour's answer is a bottom-up count on its virtual tree plus a greedy.

\textbf{Prerequisites.} \(\secref{c30-construction}\) (the build), Ch.~13
(\(\thmref{c10-subtree-template}\), the fold), Ch.~16 (\(\thmref{c11-distf}\),
distance).

\textbf{After this section you can}
\begin{itemize}[nosep]
  \item weight a compressed edge so path/cost answers are exact; \hfill\lv{5}
  \item carry the count of swallowed nodes when the DP charges per real node; \hfill\lv{5}
  \item prove a virtual-tree DP equals the honest full-tree DP (\(\thmref{c30-fold}\)); \hfill\lv{5}
  \item solve CF 2252F -- per-colour virtual tree + a count DP + a greedy. \hfill\lv{6}
\end{itemize}
\end{topicintro}

% ---------------------------------------------------------------
\subsection{What a compressed edge must remember}
% ---------------------------------------------------------------

\begin{definition}{Edge weight and skipped count on a virtual tree}{c30-edgeweight}
Let \((u,v)\) be a virtual edge with \(u\) a proper ancestor of \(v\), standing for
the real path of \(L=\depth(v)-\depth(u)\) edges.  \textbf{Weight rule:} the edge
carries \(w(u,v)=L\), so a length/sum over virtual paths equals the same over
real paths.  \textbf{Skipped-count rule:} let \(\mathit{sub}(v)\) be the number of
\emph{real} tree nodes \(x\) such that \(x\) lies on a swallowed chain below this
edge and is not itself a kept node but whose whole contribution to the DP is
\(v\)'s; concretely we carry \(\mathit{skip}(u,v)=L-1\) interior nodes if the DP is
per-node, or, for a per-terminal count, the number of marked terminals in \(v\)'s
virtual subtree \(\mathit{cnt}(v)\).  Which count to carry depends on what the DP
charges -- a distance sum charges terminals (\(\mathit{cnt}\)), a node-count
charges \(\mathit{skip}\).
\end{definition}

\begin{theorem}{A weighted DP on \(A(S)\) equals the DP on \(T\)}{c30-fold}
Let a bottom-up DP over \(T\) depend, for the marked set \(S\), only on quantities
carried along the union of paths between marked nodes (distances, subtree terminal
counts, connectedness).  Then evaluating the same recurrence on the virtual tree,
where every compressed edge contributes its weight \(L\) and its swallowed count,
returns the identical value; nodes of \(T\) that are not in any marked-to-marked path
never affect the answer and are correctly dropped.
\end{theorem}
\begin{proofbox}
Partition \(T\)'s edges into those on some path between two kept nodes (exactly the
edges covered by compressed edges of \(A(S)\)) and the rest, which hang off the
marked-spanning subtree and contain no kept node below them.  A DP whose answers for
marked nodes route all information through marked-to-marked paths is unaffected by the
second class, so dropping them changes nothing.  For the first class, a compressed
edge \((u,v)\) of length \(L\) is a chain of \(L\) real edges with no kept node in
their interior; any quantity that accumulates additively along the chain (a distance,
a sum of per-edge weights) is \(L\) copies of a constant the recurrence can multiply
in one step, and any branching count is fixed at the endpoints, which \(A(S)\) keeps.
Hence the fold over \(A(S)\) with the two carrying rules recomputes exactly the same
totals, in \(|A(S)|=O(k)\) steps instead of \(n\). \(\Box\)
\end{proofbox}

\begin{lstlisting}[caption={The generic weighted fold on a virtual tree. Same shape as Ch.13's subtree template, but each child edge contributes a per-real-edge quantity times its length, and a subtree terminal count cnt[v]. ASCII only}, label={lst:c30-fold}]
// vtKids[v] = virtual children, each (child, length L)
ll cnt[MAXN];            // marked terminals in v's virtual subtree
ll acc[MAXN];            // the folded answer
void vtDfs(int v){
    cnt[v] = marked[v] ? 1 : 0;      // v is a terminal of this colour/set
    acc[v] = 0;
    for (auto [u,L] : vtKids[v]){
        vtDfs(u);
        cnt[v] += cnt[u];
        // each of the L real edges on (v,u) carries the same flow = min(cnt[u], total-cnt[u]);
        // fold whatever per-edge quantity the problem wants:
        acc[v] += acc[u] + perEdge(cnt[u]) * L;
    }
}
\end{lstlisting}

% ---------------------------------------------------------------
\subsection{Problem Lab: CF 2252F Spectral Components -- per-colour VT + count DP + greedy}
% ---------------------------------------------------------------

\solvedproblem{c30-spectral}{Spectral Components}{CF 2252F}
  {https://codeforces.com/problemset/problem/2252/F}{CF 2600 L6}
  {per-colour virtual tree; bottom-up DP for the terminal count \(m_{\text{in}}\) on each
   virtual subtree; a virtual edge of length \(L\) has every real edge worth
   \(\min(m_{\text{in}}, m_c-m_{\text{in}})\); the connected \(k_c\)-component absorbs the
   \(k_c-1\) heaviest such edges -- a greedy over the virtual tree's edges}

\begin{problemstatement}[Spectral Components (statement and sample verified against the live CF page)]
A tree on \(n\) vertices; vertex \(i\) has colour \(c_i\in[1,n]\).  For a colour \(c\)
that appears, let \(m_c\) be its vertex count and \(k_c\le m_c\) a target size.  For
each colour \(c\) \emph{independently}, choose a connected subgraph \(S\) of exactly
\(k_c\) vertices (they need not be colour \(c\)).  The cost is \(\sum_{v:\,c_v=c}
\operatorname{dist}(v,S)\) -- the sum of shortest distances from every colour-\(c\)
vertex to \(S\).  Output, for colours \(1..n\), the minimum cost, or \(-1\) for absent
colours.  \(\sum n\le 2\cdot10^5\).

\textbf{Example (official, test 1).} \(n=5\); colours \(1\,1\,2\,1\,2\); \(k=2\,1\,1\,1\,1\);
edges \(1\!-\!2,2\!-\!3,2\!-\!4,4\!-\!5\).  Output \(\;1\;3\;-1\;-1\;-1\).
\end{problemstatement}

\begin{attemptgate}[Attempt it yourself first -- \lv{6}, give yourself 50 minutes]
Solve it for \(k_c=1\) (a single meeting vertex) first.  What does the sum of
distances to the best vertex look like when you add it up edge by edge?
\peek{root the tree; each edge with \(t\) colour-\(c\) vertices below it forces
\(\min(t, m_c-t)\) colour-\(c\) vertices to cross it unless the chosen vertex is on
the far side.  Sum that over edges to get the \(k_c=1\) baseline.}
\nudge{a connected set of \(k_c\) vertices is a connected subtree of \(k_c-1\) edges
on the paths between colour-\(c\) vertices; every edge you manage to put INSIDE \(S\)
stops charging \(\min(t,m_c-t)\).  So cost \(=\) (baseline) \(-\) (sum of the weights of
the \(k_c-1\) edges you absorb).  Minimise cost = maximise the absorbed weight.}
\fullhint{The colour-\(c\) edges that can ever matter are exactly the compressed edges
of the virtual tree \(A(S_c)\).  Do a bottom-up DP on that virtual tree: at a virtual
edge \((u,v)\) with \(v\) deeper, let \(m_{\text{in}}=\mathit{cnt}(v)\) be the colour-\(c\)
terminals below; the real edge just below \(u\) toward \(v\) has weight
\(\min(m_{\text{in}}, m_c-m_{\text{in}})\), and by the centroid-monotonicity of these
weights all \(L=\depth(v)-\depth(u)\) real edges of that compressed path share it.  The
baseline (all those edges charged) minus the heaviest \(k_c-1\) of them is the answer;
\(\sum_c m_c=n\) so all virtual trees cost \(\bigO{n\log n}\) total.}
\end{attemptgate}

\solstep{Understanding}
Two ideas stack here.  (1) \emph{Edge-contribution to a distance-to-a-set sum}: rooting
the tree, an edge whose far side holds \(t\) of the \(m_c\) colour-\(c\) vertices is
crossed by \(\min(t,m_c-t)\) of them on their way to the best component, so the total
cost is a sum of per-edge weights, and adding an edge \emph{into} the component deletes
that edge's weight.  (2) \emph{Only the virtual tree's edges have non-zero weight}:
an edge on a swallowed chain with no colour-\(c\) terminal branching below it has the
same \(t\) as its parent, so the distinct weights are exactly those of the compressed
edges of \(A(S_c)\).  Hence build the virtual tree, run a bottom-up terminal count
(the fold \(\thmref{c30-fold}\)), read each compressed edge's weight, and grab the
\(k_c-1\) heaviest.

\solstep{Key Observations}
\begin{observation}
For \(k_c=1\) the cost is \(\sum_{e} \min(t_e, m_c-t_e)\) over the edges \(e\) of the
colour-\(c\) Steiner tree; a virtual edge of length \(L\) contributes that per-edge
weight \(L\) times (\(\defref{c30-edgeweight}\)).
\end{observation}
\begin{observation}
A connected \(k_c\)-vertex component is \(k_c-1\) tree edges absorbed into \(S\); each
absorbed edge stops paying its weight.  Because the weights are non-increasing as you
move away from the tree's centroid, the \(k_c-1\) globally-heaviest weight-edges are
already connected and contain the centroid -- so a plain greedy (sort virtual edges by
weight, take \(k_c-1\)) is optimal, no DP over subsets needed.
\end{observation}

\solstep{The Approach Ladder}
\begin{center}\renewcommand{\arraystretch}{1.25}\setlength{\tabcolsep}{3pt}
\begin{tabular}{@{}c p{0.27\linewidth} p{0.31\linewidth} l c@{}}
\toprule
\textbf{\#} & \textbf{Approach} & \textbf{Idea} & \textbf{Time} & \textbf{Verdict}\\
\midrule
1 & \lv{3}\ BFS from each colour set & per colour, multi-source over the whole tree & \(\bigO{n\cdot\#\text{colours}}\) & \tle \\
2 & \lv{4}\ edge-contribution DP on the \emph{full} tree & sum \(\min(t,m_c-t)\) per colour by DFS & \(\bigO{n\cdot\#\text{colours}}\) & \tle \\
3 & \lv{5}\ same DP on the \emph{virtual} tree + greedy absorb & \(O(m_c\log m_c)\) per colour & \(\bigO{n\log n}\) & \ac \\
\bottomrule
\end{tabular}\end{center}
Rungs~1--2 pay for the whole tree once per colour and blow up when many colours are
present; rung~3 pays only for the colour's own terminals and their LCAs.

\solstep{Solution}
Root the tree once and preprocess \(\tin/\depth/\lca\).  For each colour \(c\) with
\(m_c\ge1\): build \(A(S_c)\) (\(\coderef{c30-build}\)); run the weighted fold
\(\coderef{c30-fold}\) on it to get \(\mathit{cnt}(v)\) at every kept node; for each
compressed edge \((u,v)\) compute weight \(w'=\min(\mathit{cnt}(v),m_c-\mathit{cnt}(v))\)
and multiplicity \(L=\depth(v)-\depth(u)\).  The answer is
\(\sum w'\!\cdot\! L\) minus the greedy sum of the \(k_c-1\) heaviest weight-edges
(expanding a virtual edge of multiplicity \(L\) yields \(L\) edges of weight \(w'\)).
Absent colours print \(-1\).

\solstep{Implementation}
\begin{lstlisting}[caption={CF 2252F via per-colour virtual tree + count DP + greedy absorb. For each colour build its virtual tree, fold terminal counts, each virtual edge (u,v) has weight min(cnt[v], m minus cnt[v]) repeated L = depth[v]-depth[u] times; answer = sum of weights minus the (k-1) heaviest. ASCII only}, label={lst:c30-spectral}]
#include <bits/stdc++.h>
using namespace std;
typedef long long ll; const int MAXN=200005;
int n, col[MAXN], kk[MAXN], mcnt[MAXN], tin[MAXN], dep[MAXN], tmr, par[20][MAXN];
vector<int> g[MAXN], term[MAXN];
vector<pair<int,int>> vt;                 // (node, parent) of the current colour's VT
void dfs(int v,int p,int d){ tin[v]=tmr++; dep[v]=d; par[0][v]=p;
  for(int i=1;i<20;i++) par[i][v]=par[i-1][par[i-1][v]];
  for(int u:g[v]) if(u!=p) dfs(u,v,d+1); }
int lca(int u,int v){ if(dep[u]<dep[v])swap(u,v); int d=dep[u]-dep[v];
  for(int i=0;i<20;i++) if(d>>i&1)u=par[i][u]; if(u==v)return u;
  for(int i=19;i>=0;i--) if(par[i][u]!=par[i][v]){u=par[i][u];v=par[i][v];} return par[0][u]; }
bool isAnc(int u,int v){ return tin[u]<=tin[v]; }          // within same-root DFS entry order
int cntSub[MAXN];                                          // terminals in v's VT subtree
ll solveColour(int c){
  auto &S=term[c]; int m=S.size(); if(!m) return -1;
  sort(S.begin(),S.end(),[](int a,int b){return tin[a]<tin[b];});
  vector<int> nd=S; for(int i=0;i+1<m;i++) nd.push_back(lca(S[i],S[i+1]));
  sort(nd.begin(),nd.end(),[](int a,int b){return tin[a]<tin[b];});
  nd.erase(unique(nd.begin(),nd.end()),nd.end());
  vector<int> st; st.push_back(nd[0]);
  for(int i=1;i<(int)nd.size();i++){ int v=nd[i];
    while(!isAnc(st.back(),v)) st.pop_back();
    int u=st.back();                       // virtual edge u -> v
    // weight of the real edge u-..-v below u = min(m_in, m - m_in), m_in known after fold
    vt.push_back({v,u}); st.push_back(v); }
  // bottom-up count: init cnt = 1 at terminals, 0 at pure-LCA nodes, then fold
  for(int x:nd) cntSub[x] = 0;
  for(int x:S) cntSub[x] = 1;
  sort(vt.begin(),vt.end(),[](auto&a,auto&b){return dep[a.first]>dep[b.first];});
  ll base=0; vector<pair<int,int>> wcount;  // (weight, multiplicity L)
  for(auto [v,u]:vt){ int win=cntSub[v]; cntSub[u]+=win;
    int L=dep[v]-dep[u]; int w=min(win, m-win);
    base += 1LL*w*L; for(int e=0;e<L;e++) wcount.push_back({w,1}); }
  sort(wcount.rbegin(),wcount.rend());      // heaviest weight-edges first
  ll absorb=0; int take=kk[c]-1;             // component has k_c-1 internal edges
  for(int i=0;i<(int)wcount.size() && take>0;i++){ absorb+=wcount[i].first; take--; }
  return base-absorb;
}
int main(){ int T; scanf("%d",&T); while(T--){
  scanf("%d",&n); for(int i=1;i<=n;i++){ scanf("%d",&col[i]); term[i].clear(); mcnt[i]=0;}
  for(int i=1;i<=n;i++){ scanf("%d",&kk[i]);}
  for(int i=1;i<n;i++){int u,v;scanf("%d%d",&u,&v);g[u].push_back(v);g[v].push_back(u);}
  tmr=0; dfs(1,0,0);
  for(int i=1;i<=n;i++){ term[col[i]].push_back(i); mcnt[col[i]]++; }
  for(int c=1;c<=n;c++){ vt.clear(); printf("%lld ", solveColour(c)); }
  printf("\n");
  for(int i=1;i<=n;i++) g[i].clear();
} }
\end{lstlisting}

\solstep{Trace}
\begin{tracebox}[Official test 1: two colours, root at 1]
Tree \(1\!-\!2,\,2\!-\!3,\,2\!-\!4,\,4\!-\!5\); depths \(0,1,2,2,3\); colours
\(1,1,2,1,2\); \(k=2,1,\dots\).\\
\textbf{Colour 1}, \(S=\{1,2,4\}\), \(m=3\). VT edges: \(1\!\to\!2\) (\(L=1\)) and
\(2\!\to\!4\) (\(L=1\)). Bottom-up counts: \(\mathit{cnt}(4)=1\), \(\mathit{cnt}(2)=2\).
Edge \(2\!\to\!4\): \(w=\min(1,3-1)=1\).  Edge \(1\!\to\!2\): \(w=\min(2,1)=1\).  Base
\(=1+1=2\).  \(k_1-1=1\) absorbed edge: take the heaviest (\(w=1\)) \(\Rightarrow\)
absorb \(=1\).  Cost \(=2-1=\boxed{1}\).\\
\textbf{Colour 2}, \(S=\{3,5\}\), \(m=2\). Adjacent LCA \(\lca(3,5)=2\); VT edges
\(2\!\to\!3\) (\(L=1\)), \(2\!\to\!5\) (\(L=2\), the chain \(2\text{-}4\text{-}5\)).
\(\mathit{cnt}(3)=1\Rightarrow w=1\); \(\mathit{cnt}(5)=1\Rightarrow w=\min(1,1)=1\) over
\(L=2\) real edges.  Base \(=1+2=3\).  \(k_2-1=0\) absorbed \(\Rightarrow\) cost
\(\boxed{3}\).\\
Colours \(3,4,5\) absent \(\Rightarrow -1\).  Output \(1\;3\;-1\;-1\;-1\) -- matching the
official sample.  (Tests 2 and 3 give \(3\,0\,-1\ldots\) and \(3\,2\,-1\ldots\);
the same fold on their virtual trees reproduces them.)
\end{tracebox}

\begin{complexitybox}[CF 2252F]
Per colour the virtual tree is \(\bigO{m_c\log m_c}\) (a sort plus \(m_c\) adjacent LCAs)
and the fold is \(\bigO{m_c}\); since \(\sum_c m_c=n\), the total is \(\bigO{n\log n}\)
after the \(\bigO{n\log n}\) LCA preprocess, versus \(\bigO{n\cdot\#\text{colours}}\) for a
full-tree DP per colour.  Space \(\bigO{n}\).
\end{complexitybox}

\begin{warningbox}[Where the weighted fold goes wrong]
\begin{enumerate}[nosep,leftmargin=1.6em]
  \item \textbf{Dropping the \(L\) factor.} A compressed edge is \(L\) real edges of the
        \emph{same} weight; base and absorb both count edges, so a virtual edge of
        length \(L\) contributes \(L\) selectable weight-\(w'\) edges.  Forgetting \(L\)
        under-costs every deep chain.
  \item \textbf{Wrong \(\mathit{cnt}\).} The bottom-up count must run on the virtual tree
        (terminals start at \(1\), pure-LCA nodes at \(0\)); counting real subtree sizes
        instead gives the wrong \(w\).
  \item \textbf{Absorbing more than \(k_c-1\) edges} makes \(S\) too large; the greedy
        takes exactly \(k_c-1\) and stops.
\end{enumerate}
\end{warningbox}

\begin{hardeningbox}
\begin{itemize}[nosep,leftmargin=1.4em]
  \item \textbf{Greedy-uses-monotonicity.} The \(\min(t,m-t)\) weights are non-increasing
        away from the centroid, so the \(k_c-1\) heaviest edges are connected and legal --
        a greedy is provably optimal; do not reach for a subset DP.
  \item \textbf{Rebuild \(\mathit{vt}\)/counts per colour} and clear; stale virtual
        children across colours corrupt the fold.
  \item \textbf{\(-1\) for absent colours} is emitted by \(m=0\), not by an empty VT.
  \item \textbf{Use \texttt{long long}} for base and absorb; costs reach \(\sim n^2/4\).
\end{itemize}
\end{hardeningbox}

\begin{recallbox}[26.2 in one breath]
\emph{A virtual tree is faithful once every compressed edge carries its length
\(\depth(v)-\depth(u)\) and the count the DP charges (terminals \(\mathit{cnt}\) or
skipped \(\mathit{skip}\)); then Ch.~13's bottom-up fold on \(A(S)\) equals the fold on
\(T\) while touching only \(O(k)\) nodes (\(\thmref{c30-fold}\)).  CF 2252F is that
machine: per colour, build \(A(S_c)\), fold the terminal counts so a virtual edge is
worth \(\min(m_{\text{in}},m_c-m_{\text{in}})\) repeated \(L\) times, then subtract the
\(k_c-1\) heaviest weight-edges -- the centroid-monotonicity makes the greedy exact.}
\selftest{
  \item Why can nodes outside every marked-to-marked path be dropped? \emph{They hold no kept node below and a distance/count DP routes nothing through them. \(\thmref{c30-fold}\).}
  \item What two quantities must a compressed edge remember? \emph{Its real length \(L\) and the count the DP charges on it (\(\defref{c30-edgeweight}\)).}
  \item In 2252F, why is absorbing the heaviest \(k_c-1\) edges optimal? \emph{Weights are centroid-monotone, so the heaviest are already connected. Observation~2.}
}
\end{recallbox}
